% !TEX root = main.tex
% Method section for the Attack4dLLM paper.
% Required preamble:
%   \usepackage{algorithm}
%   \usepackage{algpseudocode}
%   \usepackage{amsmath,amssymb,bm}
%   \newcommand{\method}{\textsc{Search}\xspace}
%   \newcommand{\dual}{\textsc{Dual100}\xspace}
%   \newcommand{\maskTok}{\langle\textsc{mask}\!:\!N\rangle}

\section{Method}
\label{sec:method}

\subsection{Problem setup and notation}
\label{sec:method:setup}

Let $g$ denote a harmful \emph{goal} drawn from a benchmark
$\mathcal{G}=\{g_1,\dots,g_{|\mathcal{G}|}\}$ (HarmBench, JailbreakBench, or
StrongReject). Let $V_\theta$ denote a diffusion LLM victim that, given a
prompt $p$ and a mask-bearing template $\tau$ containing $K$ masked spans,
produces a completion $V_\theta(p,\tau)\in\Sigma^*$ by parallel iterative
denoising. A \emph{template} is a string of the form
$\tau = u_0\,\maskTok_{\!1}\,u_1\,\maskTok_{\!2}\,\dots\,\maskTok_{\!K}\,u_K$
where each $u_j$ is non-mask text and each $\maskTok_{\!j}$ is later expanded
to a victim-specific number of native mask tokens (e.g. $\langle\,$\verb|mask|$\,\rangle$
for \llada\,2 / Dream, $\langle\,$\verb|mdm_mask|$\,\rangle$ for
\llada\,/\llada\,1.5\,/\mmada). A \emph{judge} $J$ assigns a binary success
label and a $0$-$5$ harmscore to the pair $(g,V_\theta(p,\tau))$.

A \emph{pattern} $\pi$ is a structural schema
$\pi=(\textsf{type},\textsf{org},K,\bm{r},\mathcal{T})$ specifying the
template's structure type (e.g.\ procedural-steps, narrative-arc, worked-example),
organization style, slot count $K$, slot-role vector
$\bm{r}=(r_1,\dots,r_K)$, and a set $\mathcal{T}$ of goal-type tags
indicating which categories of $g$ this pattern is suited to.
$\Pi_t$ denotes the \emph{pattern registry} (library) at iteration $t$;
$\Pi_0$ is initialized in Stage~A (Section~\ref{sec:method:exploration}) and
grows monotonically as new patterns are distilled (Section~\ref{sec:method:evolution}).

\subsection{Architecture overview}
\label{sec:method:overview}

The pipeline (Fig.~\ref{fig:pipeline}, Algorithm~\ref{alg:expansion}) is a
two-stage system.
\textbf{Stage~A (offline, run once).} A small bootstrap set of harmful goals is
fed to the attacker LLM with a generic ``produce a mask-bearing template''
instruction; the resulting templates are abstracted into pattern schemas,
deduplicated by stable hashing, and written to the shared registry on disk.
\textbf{Stage~B (online, the core loop).} For each new goal $g$, a UCB bandit
selects a pattern from $\Pi_t$, the attacker instantiates a concrete template,
the dLLM victim fills the masks, and the scorer assigns a reward. Low-reward
attempts trigger a base-model fallback; high-reward attempts that improve on
the previous iteration trigger pattern distillation. Both events update
$\Pi_t\!\to\!\Pi_{t+1}$, so attack effort spent on goal $g_i$ benefits goal
$g_{i+1}$.

We instantiate four roles with three distinct LLMs:
the \emph{attacker} and \emph{scorer} use Qwen3-4B-Instruct served via vLLM;
the \emph{weaker base model} (used only inside the fallback) is the
unaligned Qwen3-4B-Base; the optional \emph{tagger} for the fallback's
redaction step is Qwen3-235B served via AWS Bedrock (with a local
Qwen3-4B-Instruct fallback when Bedrock is unavailable).

\subsection{Stage A: exploration}
\label{sec:method:exploration}

Given a bootstrap set $\mathcal{G}_0\subset\mathcal{G}$ of $|\mathcal{G}_0|=150$
goals, we elicit one mask-bearing template per goal from the attacker.
Each template is run through an \emph{extractor} (a single LLM call) that
returns the schema $\pi$. Patterns are keyed by
$\textsf{id}(\pi)=\textsc{Hash}(\pi)$ so semantically equivalent
schemas collapse to the same registry entry. The result is a starting
library $\Pi_0$ of $\sim 30$--$50$ unique patterns, persisted to
\verb|logs/shared_pattern_registry.json|.

\subsection{Stage B: UCB-guided online search}
\label{sec:method:expansion}

For each goal $g$ we run $T=\mathtt{max\_iters}$ iterations
(Algorithm~\ref{alg:expansion}). At iteration $t$ we:
\begin{enumerate}
\itemsep0pt
\item \textbf{Filter the registry by goal type.} A retriever (one LLM call)
      maps $g$ to one or more goal-type tags
      $\mathcal{T}(g)\subseteq\{\text{cybercrime,harassment,\ldots}\}$ and we
      keep $\Pi_t(g)=\{\pi\in\Pi_t:\mathcal{T}(g)\cap\mathcal{T}_\pi\neq\emptyset\}$.
\item \textbf{UCB selection.} For each $\pi\in\Pi_t(g)$ we maintain
      $(\hat\mu_\pi,n_\pi)$ (running mean reward, visit count). Unvisited
      patterns are selected first; otherwise we pick
      \begin{equation}
      \pi^\star_t \;=\; \argmax_{\pi\in\Pi_t(g)}\;
        \hat\mu_\pi + \alpha \sqrt{\frac{2\ln t}{n_\pi}},
      \qquad \alpha=1.0.
      \label{eq:ucb}
      \end{equation}
\item \textbf{Attacker.instantiate.} A \emph{single} LLM call to the attacker,
      conditioned on $g$ and $\pi^\star_t$, emits a four-part response:
      a strategy tag, a jailbreak prompt $p$, and a $K_\pi$-slot template
      $\tau$ (with one continuation example used only for self-supervision).
      We post-process $\tau$ via \textsc{RescaleMasks} so that the total
      number of native mask tokens hits a target ratio of the victim's
      \texttt{max\_gen\_length}: $\sim 13$ tokens$\times 7$ slots
      ($\approx\!36\%$) for \method, $\sim 55$ tokens$\times 5$ slots
      ($\approx\!64\%$) for \dual.
\item \textbf{Victim infill.} We dispatch to the appropriate
      \texttt{sample/} backend by victim, expand $\maskTok$ into the
      victim's native mask token, and decode in parallel for
      \texttt{steps} denoising iterations to obtain
      $y_t = V_\theta(p,\tau)$.
\item \textbf{Scorer.} A single LLM call returns
      $r_t\in\{0,1\}$ along with a short evidence string and an
      analysis used for fallback feedback.
\item \textbf{Fallback (only if $r_t<\rho$).} With $\rho=0.7$ and at most
      $R=5$ retries we run Algorithm~\ref{alg:fallback}; if any retry
      improves on the current $r_t$ we adopt that $(\tau,y)$ pair.
\item \textbf{Bandit update.} $n_{\pi^\star_t}\!\leftarrow n_{\pi^\star_t}+1$;
      $\hat\mu_{\pi^\star_t}\!\leftarrow
       \hat\mu_{\pi^\star_t}+(r_t-\hat\mu_{\pi^\star_t})/n_{\pi^\star_t}$.
\item \textbf{Pattern evolution.} If $t\!\ge\!2$ and $r_t\!\ge\!r_{t-1}$ we run
      the summarizer (Section~\ref{sec:method:evolution}); the resulting
      schema--template pair is added to $\Pi_{t+1}$ if its hash is new.
\end{enumerate}

\begin{algorithm}[t]
\caption{\method: per-goal UCB attack loop (Stage~B).}
\label{alg:expansion}
\begin{algorithmic}[1]
\Require goal $g$, registry $\Pi_t$, attacker $A$, victim $V_\theta$, scorer $S$,
         summarizer $\Sigma$, extractor $E$, hyperparameters $T,\alpha,\rho,R$
\State $\Pi_t(g)\leftarrow\textsc{FilterByGoalType}(\Pi_t,g)$
\State $\theta\leftarrow\{(\hat\mu_\pi=0,n_\pi=0):\pi\in\Pi_t(g)\}$
\State $(\pi_{\text{prev}},\tau_{\text{prev}},r_{\text{prev}})\leftarrow(\bot,\bot,\bot)$
\For{$t=1,\dots,T$}
  \State $\pi^\star\leftarrow\textsc{SelectUCB}(\Pi_t(g),\theta,t,\alpha)$ \Comment{Eq.~\ref{eq:ucb}}
  \State $(p,\tau)\leftarrow A.\textsc{Instantiate}(g,\pi^\star)$
  \State $y\leftarrow V_\theta(p,\tau)$
  \State $r\leftarrow S.\textsc{Score}(g,p,\tau,y)$
  \If{$r<\rho$}
    \State $(\tau,y,r)\leftarrow\textsc{Fallback}(g,\pi^\star,r,R)$ \Comment{Algorithm~\ref{alg:fallback}}
  \EndIf
  \State $\textsc{UpdateUCB}(\theta,\pi^\star,r)$
  \If{$t\ge 2 \;\land\; r\ge r_{\text{prev}}$}
    \State $(\tilde\tau,\tilde\pi)\leftarrow\Sigma.\textsc{Distill}(\pi_{\text{prev}},\tau_{\text{prev}},r_{\text{prev}},\pi^\star,\tau,r)$
    \If{$\tilde\pi$ has empty schema} $\tilde\pi\leftarrow E.\textsc{Extract}(\tilde\tau)$ \EndIf
    \If{$\textsc{Hash}(\tilde\pi)\notin\Pi_t$} $\Pi_t\leftarrow\Pi_t\cup\{(\tilde\pi,\tilde\tau)\}$ \EndIf
  \EndIf
  \State $(\pi_{\text{prev}},\tau_{\text{prev}},r_{\text{prev}})\leftarrow(\pi^\star,\tau,r)$
\EndFor
\State \Return $\arg\max_{(\tau,y,r)\text{ seen}} r$
\end{algorithmic}
\end{algorithm}

\subsection{Scorer-guided fallback}
\label{sec:method:fallback}

The fallback (Algorithm~\ref{alg:fallback}) is the mechanism by which the
attack escapes pattern-level dead ends. Its key idea is that an unaligned
\emph{base} model can fluently write content that an aligned attacker cannot,
but base outputs are not directly usable as templates because they contain
no masks. We therefore (i) draft a raw response, (ii) tag the
goal-realizing spans with a stronger model, and (iii) convert tags to masks
to yield a goal-specific template that the dLLM victim is then asked to
re-fill itself. Concretely:

\begin{algorithm}[t]
\caption{\textsc{Fallback}: scorer-guided template repair.}
\label{alg:fallback}
\begin{algorithmic}[1]
\Require goal $g$, current pattern $\pi$, current best reward $r$, retries $R$
\State $\textit{feedback}\leftarrow\text{scorer's last analysis (truncated to 600 chars)}$
\For{$i=1,\dots,R$}
  \State $\tilde g\leftarrow A.\textsc{RephraseGoal}(g,\textit{framing}_i)$
         \Comment{rotates over: direct, narrative, tutorial, imperative, worked-example}
  \State $z\leftarrow A_{\text{base}}.\textsc{Continue}(\tilde g)$
         \Comment{weaker, unaligned base model}
  \State $z^{\text{tagged}}\leftarrow A_{235\text{B}}.\textsc{TagSpans}(g,z)$
         \Comment{wraps goal-realizing spans in $[\textsc{redact}]\cdots[/\textsc{redact}]$}
  \State $\tau'\leftarrow\textsc{TagsToMasks}(z^{\text{tagged}})$
  \State $y'\leftarrow V_\theta(p,\tau')$;\quad
         $r'\leftarrow S.\textsc{Score}(g,p,\tau',y')$
  \If{$r'>r$} $(\tau,y,r)\leftarrow(\tau',y',r')$;\;\;
         $\textit{feedback}\leftarrow$ scorer's new analysis
  \EndIf
  \If{$r\ge\rho$} \textbf{break} \EndIf
\EndFor
\State \Return $(\tau,y,r)$
\end{algorithmic}
\end{algorithm}

Two design points are worth flagging.
\textbf{(F1) Goal rephrasing rotates framings.} Calling the base model on the
literal goal often hits its residual safety training; rotating the framing
(narrative, tutorial, imperative, worked-example) produces stylistically
distinct continuations and gives the bandit a richer set of candidate
templates without the attacker LLM having to ``know'' which framing the
victim will accept.
\textbf{(F2) Tagging happens with a stronger model than drafting.} The
draft model needs fluency, not judgement; the tagger needs judgement, not
generation breadth. Splitting the two roles and routing the second to
Qwen3-235B (or a local instruct fallback) yields cleaner mask placement
than asking a single model to do both.

\subsection{Pattern evolution (summarizer)}
\label{sec:method:evolution}

When iteration $t$ improves on iteration $t-1$ for the same goal,
the summarizer is given both $(\pi_{t-1},\tau_{t-1},r_{t-1})$ and
$(\pi_t,\tau_t,r_t)$ and asked, in a single LLM call, to produce a
\emph{new} schema--template pair $(\tilde\pi,\tilde\tau)$ that captures
what made the second attempt better. If the returned schema is malformed
we re-extract from $\tilde\tau$ via the extractor used in Stage~A. The
pair is added to $\Pi_t$ keyed by $\textsc{Hash}(\tilde\pi)$, so duplicates
collapse. Patterns therefore have three provenances:
\textsc{exploration} (Stage~A), \textsc{expansion} (Stage~B distillation),
and \textsc{fallback} (templates promoted out of Algorithm~\ref{alg:fallback}
when they win their iteration). The registry is shared across goals and
benchmarks but partitioned by victim, so transfer happens \emph{within} a
victim only.

\subsection{The \dual variant}
\label{sec:method:dual100}

\dual differs from \method only in two hyperparameters. (1) Mask granularity:
\dual emits $K{=}5$ paragraph-sized slots ($\sim 55$ native mask tokens each)
versus \method's $K{=}7$ entity-sized slots ($\sim 13$ tokens each). (2)
Attacker contribution: because \dual's slots are larger, the victim
generates $\approx 64\%$ of the final tokens versus $\approx 36\%$ for
\method. The two variants share Algorithm~\ref{alg:expansion},
Algorithm~\ref{alg:fallback}, the registry, and all other hyperparameters.
\dual is the variant intended for collecting safety-alignment training
data, since the victim's contribution dominates the final output and is
therefore informative as a model-specific failure mode; \method is the
variant intended for adversarial robustness benchmarking, where having the
attacker write more of the surface text yields a stricter test of whether
the victim's safety alignment generalizes off-distribution.

\subsection{Hyperparameters and compute}
\label{sec:method:hparams}

Unless stated otherwise we use $T=5$, $\alpha=1.0$, $\rho=0.7$, $R=5$,
$|\mathcal{G}_0|=150$, attacker temperature $0.7$, scorer temperature $0.0$,
and victim \texttt{steps} matching each model's released default
($256$ for \llada\,2, $128$ for \llada/\llada\,1.5/\mmada, $256$ for Dream).
A full run on $100$ goals against one victim costs roughly
$100\!\times\!T\!\times\!(1\!+\!\bar R)$ attacker LLM calls (where
$\bar R\!\approx\!1.4$ is the empirical average fallback retry count) plus
the same number of victim denoising passes; on a single H100 this is
$\sim 6$--$9$ wall-clock hours per (victim, benchmark) pair.
